\documentclass[11pt,a4paper]{article}
\usepackage[utf8]{inputenc}
\usepackage[T1]{fontenc}
\usepackage[margin=2.4cm]{geometry}
\usepackage{graphicx}
\graphicspath{{../img/}}
\usepackage{booktabs}
\usepackage{longtable}
\usepackage{array}
\usepackage{enumitem}
\usepackage{float}
\usepackage{caption}
\usepackage{fancyhdr}
\usepackage[hidelinks]{hyperref}

\captionsetup{font=small,labelfont=bf}
\setlist{nosep,leftmargin=1.4em}
\pagestyle{fancy}\fancyhf{}
\lhead{\small ScholarPath --- Class Diagrams}\rhead{\small\thepage}
\renewcommand{\headrulewidth}{0.4pt}

\newcommand{\diagram}[2]{%
  \begin{figure}[H]\centering
  \includegraphics[width=\textwidth,height=0.66\textheight,keepaspectratio]{#1}%
  \caption{#2}\end{figure}}

\title{\textbf{ScholarPath --- Class Diagrams (UML)}\\[2pt]
\large Domain model and architecture (Sommerville notation)}
\author{ScholarPath Graduation Project}
\date{\today}

\begin{document}
\maketitle\thispagestyle{fancy}

\begin{abstract}\noindent
This report presents the ScholarPath domain model as UML class diagrams in the
style of Sommerville (\emph{Software Engineering}, 9th ed.). To keep every figure
legible, the model is split by subject area; a separate diagram shows the base
type hierarchy, and a final ``ports \& adapters'' diagram captures the Clean
Architecture seam between the application and infrastructure layers. Classes,
attributes and key operations are taken directly from
\texttt{ScholarPath.Domain/Entities} and the application interfaces.
\end{abstract}

\tableofcontents\bigskip

\section{Notation}
\begin{itemize}
  \item A \textbf{class} is a box with name / attributes / operations.
        \texttt{<<abstract>>} marks an abstract class; \texttt{<<interface>>} an interface.
  \item \textbf{Generalization} (inheritance) is a hollow triangle pointing at the
        parent ($\triangleleft$, drawn \texttt{Base <|-- Derived}).
  \item \textbf{Realization} (a class implements an interface) is a dashed hollow
        triangle (\texttt{Interface <|.. Class}).
  \item \textbf{Composition} (filled diamond) means the part cannot exist without
        the whole; \textbf{aggregation} (hollow diamond) means it can.
  \item \textbf{Associations} carry role names and multiplicities
        ($1$, $0..1$, $*$).
\end{itemize}

\section{Base entity types}
Almost every persistent class inherits \texttt{AuditableEntity} (itself a
\texttt{BaseEntity}); classes that support soft delete also implement
\texttt{ISoftDeletable}. To avoid drawing this inheritance on every figure, the
hierarchy is shown once here and implied thereafter.
\diagram{Class_Base_Types.pdf}{Base entity types and the domain-event mechanism.}

\section{Identity \& Profile}
\diagram{Class_Identity_Profile.pdf}{Identity \& Profile classes.}
\noindent\textbf{Inheritance note.} Unlike every other persistent class,
\texttt{ApplicationUser} does \emph{not} inherit \texttt{AuditableEntity}: it
extends \texttt{IdentityUser<Guid>} and implements \texttt{ISoftDeletable}, and
therefore carries only a \emph{stripped} audit triple
(\texttt{CreatedAt, UpdatedAt, RowVersion}) --- it has no
\texttt{CreatedByUserId}\,/\,\texttt{UpdatedByUserId}. \texttt{ApplicationRole}
extends \texttt{IdentityRole<Guid>}.
\par\noindent\texttt{ApplicationUser} owns a single \texttt{UserProfile} (which
composes \texttt{EducationEntry} rows), holds many \texttt{ApplicationRole}s via
the \texttt{UserRoles} junction, issues \texttt{RefreshToken}/
\texttt{PasswordResetToken} rows, and may raise \texttt{UpgradeRequest}s
(which compose \texttt{UpgradeRequestFile}/\texttt{UpgradeRequestLink}).
\texttt{ExpertiseTag} is a lookup table. \texttt{/FullName} is a derived property.

\section{Scholarships, Applications \& Documents}
\diagram{Class_Scholarships_Applications.pdf}{Scholarships, Applications \& Documents classes.}
\noindent A \texttt{Scholarship} (optionally classified by a \texttt{Category})
receives many \texttt{ApplicationTracker}s. An application can be backed by a
\texttt{Document} and finalised by exactly one \texttt{CompanyReview}.
\texttt{CompanyReviewRequest} models the paid-support engagement.
\texttt{ApplicationTracker.IsActive} / \texttt{IsReadOnly} are computed guards that
encode the application state machine.

\section{Booking, Payments \& Ratings}
\diagram{Class_Booking_Payments_Ratings.pdf}{Booking, Payments \& Ratings classes.}
\noindent A \texttt{ConsultantBooking} reserves a \texttt{ConsultantAvailability},
is settled by one \texttt{Payment}, may be rated by one \texttt{ConsultantReview},
and is composed of \texttt{SessionRecording}s. A \texttt{Payout} aggregates many
\texttt{Payment}s. Monetary amounts are integer \emph{cents}.

\section{Community \& Chat}
\diagram{Class_Community_Chat.pdf}{Community \& Chat classes.}
\noindent \texttt{ForumPost} is self-referencing (threaded replies), composes
\texttt{ForumPostAttachment}s, receives \texttt{ForumVote}/\texttt{ForumFlag}/
\texttt{ForumBookmark} rows, and is tagged many-to-many with \texttt{ForumTag} via
\texttt{ForumPostTag}. \texttt{ChatConversation} composes \texttt{ChatMessage}s;
\texttt{UserBlock} records a directed block between two users.

\section{Resources Hub \& Notifications}
\diagram{Class_Resources_Notifications.pdf}{Resources Hub \& Notifications classes.}
\noindent A \texttt{Resource} is split into \texttt{ResourceChild} chapters and is
saved/tracked through \texttt{ResourceBookmark} and \texttt{ResourceProgress}
(itself composing \texttt{ResourceProgressChild}). \texttt{Notification} delivery
preferences live in \texttt{NotificationPreference}.

\section{AI, Knowledge, Platform \& Cross-cutting}
\diagram{Class_AI_Platform.pdf}{AI, Knowledge, Platform \& Cross-cutting classes.}
\noindent \texttt{AiInteraction} may be sampled by an
\texttt{AiRedactionAuditSample} and attributed to many
\texttt{RecommendationClickEvent}s; \texttt{KnowledgeDocument} backs the RAG
pipeline. \texttt{AuditLog}, \texttt{UserDataRequest}, \texttt{UserRiskFlag},
\texttt{SuccessStory} and \texttt{PlatformSetting} are the cross-cutting tables;
those marked \texttt{: BaseEntity} are intentionally \emph{not} auditable
(append-only / analytics rows).

\section{Architecture --- ports \& adapters}
The application layer defines \textbf{ports} (interfaces); the infrastructure
layer provides \textbf{adapters} (implementations), chosen at start-up by
dependency injection / configuration. This seam is why the platform can swap a
development stub for a real Azure service without touching application code.
\diagram{ScholarPath_Ports_Adapters.pdf}{Ports (interfaces) and their infrastructure adapters.}

\begin{longtable}{@{}p{4.3cm}p{6.8cm}p{3.4cm}@{}}
\toprule
\textbf{Port (interface)} & \textbf{Production adapter} & \textbf{External system}\\
\midrule\endhead
\texttt{IApplicationDbContext} & \texttt{ApplicationDbContext} (EF Core) & Azure SQL\\
\texttt{ITokenService} & \texttt{TokenService} (JWT signer) & ---\\
\texttt{IPasswordHasher} & \texttt{IdentityPasswordHasher} & ---\\
\texttt{ISsoService} & \texttt{SsoService} / \emph{StubSsoService} & Google + Microsoft OAuth\\
\texttt{IEmailVerificationService} & \texttt{EmailVerificationService} & ---\\
\texttt{IEmailChangeService} & \texttt{EmailChangeService} & ---\\
\texttt{IStripeService} & \texttt{StripeService} / \emph{StubStripeService} & Stripe\\
\texttt{IMeetingService} & \texttt{AzureCommunicationMeetingService} / \emph{StubMeetingService} & Azure Comm. Svcs\\
\texttt{IEmailService} & \texttt{MailKitEmailService} / \emph{StubEmailService} & SMTP\\
\texttt{IBlobStorageService} & \texttt{FileStorageService} & Azure Blob / local\\
\texttt{IFileScanService} & \texttt{ClamAvFileScanService} / \emph{NoOpFileScanService} & ClamAV\\
\texttt{IAiService} & \texttt{AzureOpenAiService} / \texttt{OpenAiService} / \emph{LocalAiService} & Azure OpenAI \emph{or} OpenAI direct\\
\texttt{IEmbeddingService} & \texttt{AzureOpenAiEmbeddingService} / \texttt{OpenAiEmbeddingService} / \emph{LocalEmbeddingService} & Azure OpenAI \emph{or} OpenAI direct\\
\texttt{IKnowledgeRetriever} & \texttt{KnowledgeRetriever} & (RAG read)\\
\texttt{IKnowledgeBaseIndexer} & \texttt{KnowledgeBaseIndexer} & (RAG write)\\
\texttt{IDatasetProvider} & \texttt{EmbeddedDatasetProvider} & (bundled corpora)\\
\texttt{IFieldEncryptionService} & \texttt{AesGcmFieldEncryptionService} & ---\\
\texttt{IFieldEncryptionKeyProvider} & \texttt{KeyVault\dots} / \emph{Local\dots} & Azure Key Vault\\
\texttt{IJwtKeyProvider} & \texttt{KeyVaultJwtKeyProvider} / \emph{Local\dots} & Azure Key Vault\\
\texttt{IEventPublisher} & \texttt{EventHubPublisher} / \emph{StubEventPublisher} & Azure Event Hub\\
\texttt{IPowerBiService} & \texttt{PowerBiService} / \emph{StubPowerBiService} & Power BI\\
\texttt{INotificationDispatcher} & \texttt{NotificationDispatcher} & in-app + email + SignalR\\
\texttt{IChatRealtimeNotifier} & \texttt{ChatRealtimeNotifier} & SignalR ChatHub\\
\texttt{ICommunityRealtimeNotifier} & \texttt{CommunityRealtimeNotifier} & SignalR CommunityHub\\
\texttt{IChatPresenceQuery} & \texttt{PresenceTracker} & (in-memory)\\
\texttt{IAuditService} & \texttt{AuditService} & ---\\
\texttt{IUserAdministration} & \texttt{UserAdministration} & ---\\
\texttt{IAdminReadService} & \texttt{AdminReadService} & ---\\
\texttt{IConsultantReadService} & \texttt{ConsultantReadService} & ---\\
\texttt{IChatContactReadService} & \texttt{ChatContactReadService} & ---\\
\texttt{ICurrentUserService} & \texttt{CurrentUserService} & (HttpContext)\\
\texttt{IDateTimeService} & \texttt{DateTimeService} & ---\\
\textit{13 job ports} & \textit{13 job classes (Infrastructure/Jobs)} & Hangfire (when enabled)\\
\bottomrule
\end{longtable}
\noindent Adapters shown in \emph{italics} are the development / configuration
fall-backs (\texttt{Stub*}, \texttt{NoOp*}, \texttt{Local*}) selected when the
corresponding cloud service is not configured. The AI and embedding ports each
have \emph{three} adapters: Azure OpenAI, OpenAI-direct, and an offline local
implementation, chosen by the \texttt{Ai:Provider} setting.
\end{document}
